\chapter{Full-Profile Limitations (and Why They Rarely Bite)}

The \texttt{full} profile is a complete GNARL tasking runtime rebuilt on bare
metal: no Ravenscar restrictions, with rendezvous, protected objects with
entries, dynamic and nested tasks, dynamic priorities, task attributes,
exception propagation, \texttt{abort}, and \texttt{Ada.Interrupts} handlers all
running on hardware. A handful of rough edges remain. This chapter lists them
honestly and, for each, explains why it seldom matters for a real embedded
program. The short version: the gaps cluster at the margins of the language
(constructs a startup-time, fixed-task-set application almost never reaches), and
each has a clean workaround or is explicitly permitted by the Reference Manual.

\section{Toolchain-bound: asynchronous \texttt{select} with a relative delay}\index{asynchronous transfer of control}\index{select statement}

A \texttt{select \ldots\ delay D; then abort \ldots\ end select} (the relative,
\texttt{Duration}-triggered form of asynchronous transfer of control, ACATS
\texttt{C974*}) does not compile under \texttt{full}. This is \emph{not} a runtime
gap: under \texttt{Configurable\_Run\_Time} the GNAT front end emits a call that
passes a \texttt{Duration} to \texttt{Enqueue\_RT (T : Real\_Time.Time)}, which,
crossed with this port's \texttt{type Time}, does not type-check. Fixing it is a
compiler change or a deep \texttt{Real\_Time.Time} redesign, not an RTS unit we
can add.

\textbf{Why it rarely bites.} The \emph{absolute} form (\texttt{select \ldots\
delay until T}) and ATC triggered by an \emph{entry call} both work. A relative
timeout is sugar for ``compute \texttt{Clock + D}, then wait until that instant'',
which the absolute form expresses directly. Timed operations are also commonly
written with a guard plus \texttt{Ada.Real\_Time} arithmetic, untouched by this.

\section{RM-permitted: aborting a pure CPU loop}\index{abort statement}

\texttt{abort} of a task blocked at any \emph{abort-completion point}
(a \texttt{delay}, an entry call, a rendezvous) works and is HW-verified
(\texttt{'Terminated} flips, the task unwinds cleanly). The one case that does
\emph{not} take effect is a task spinning in a tight computational loop with no
completion point at all: it keeps running until it next reaches one.

\textbf{Why it rarely bites.} This is explicitly allowed: ARM 9.8 only
requires abort to take effect at completion points, and a loop with no such point
is permitted to run on. Delivering it anyway would mean injecting an exception
into a preempted task from the timer ISR, which (as the context-switch and
interrupts chapters describe) is deep, windowed-Xtensa, shared-kernel surgery
for a behavior the standard does not mandate. In practice any loop that
does I/O, delays, or calls an entry is abortable; a genuinely tight numeric
kernel that must be cancellable can poll a \texttt{Volatile} flag.

\section{\texttt{Configurable\_Run\_Time} finalization}\index{finalization}\index{Configurable\_Run\_Time}\index{Ada.Interrupts}

The front end does not generate \emph{library-level} finalization because
\texttt{system.ads} sets \texttt{Configurable\_Run\_Time => True}. Two consequences:

\begin{itemize}
  \item A \emph{library-level} controlled object's \texttt{Finalize} runs only if
        the object is finalized at scope exit, not as part of program shutdown.
  \item Consequently, the \emph{implicit} restore-on-finalization of a
        \emph{library-level} interrupt-handler protected object does not run
        automatically at shutdown. (Both the static \texttt{pragma Attach\_Handler}
        path \emph{and} the full dynamic \texttt{Ada.Interrupts} operations
        (\texttt{Attach\_Handler}, \texttt{Detach\_Handler}, \texttt{Exchange\_-
        Handler}, \texttt{Current\_Handler}) are now supported and
        HW-verified; see the Interrupts chapter. Only this one implicit-restore
        case is outstanding.)
\end{itemize}

\textbf{Why it rarely bites.} Scope-level finalization -- the overwhelmingly
common case (a controlled object in a subprogram or block) -- works normally.
Library-level objects live for the entire life of the program anyway, so on an
embedded target that never exits there is nothing to finalize. And interrupt
handlers are wired up once at elaboration and stay for the life of the system, so
the missing implicit restore-at-shutdown never matters in practice; explicit
\texttt{Detach\_Handler} works if you do need it.

\section{\texttt{Ada.Calendar} resolution: 5\,ns, with the full range to 2399}\index{Ada.Calendar}

The bare \texttt{Ada.Calendar} originally declared \texttt{Year\_Number} as
\texttt{1901\,..\,2099} (the Ada~95 bound), whereas Ada~2005 onward (and the
modern ACATS \texttt{C96004}) require \texttt{2399}. The obstacle was the time
representation: the Calendar uses a \emph{Julian-day fixed-point} base sized so the
span fits in signed 64 bits at sub-second precision, and stretching the range to
2399 (roughly five centuries) overflows that base at 1\,ns granularity, which is
why mainline GNAT eventually abandoned the Julian-day Calendar for an
integer-nanosecond one.

Rather than swap in a new time base (a risky rewrite of a unit shared by the
\texttt{embedded} and \texttt{full} profiles), the range is bought by
\emph{coarsening} the existing base by exactly $5\times$: the sub-day fixed-point
scale changes from \texttt{Duration'Small\,/\,86\_400.0} to
\texttt{/\,17\_280.0}\,($=86\_400/5$), and the Modified-Julian-Day and Julian-Day
type ranges widen to cover 2399-12-31 (\texttt{Time'Last} compile-checks to
\texttt{197\_641}). The 1858 epoch and the day-arithmetic algorithm are untouched.
The trade is precision for range: the Calendar's resolution becomes $5$\,ns instead
of $1$\,ns, and in exchange \texttt{Year\_Number} now runs the full
\texttt{1901\,..\,2399}. \texttt{C96004} \textbf{PASSES} (HW-verified, both
profiles), and the change is applied by \texttt{gen\_runtime.sh} when it
materialises the runtime.

\textbf{Why the 5\,ns rarely bites.} \texttt{Ada.Calendar} is a wall-clock: its job
is civil date-and-time stamping, where a five-nanosecond quantum on a value that
counts whole seconds is meaningless. Any code that genuinely needs sub-microsecond
timing uses \texttt{Ada.Real\_Time} instead, whose \texttt{Time\_Span} keeps its
full \texttt{Duration'Small} (1\,ns) precision and is untouched by this change.
Arithmetic, formatting, \texttt{Split}/\texttt{Time\_Of}, and \texttt{delay until}
on a \texttt{Calendar.Time} remain HW-verified across the whole 1901--2399 span.

\section{Per-task heap reclamation under task churn}\index{task churn}

A terminated task's primary stack is returned to the heap, but a program that
\emph{churns} tasks (creating and destroying them in a tight loop) can still
exhaust the heap: the reclamation of the per-task resources is not yet fully
race-free against the terminating thread, and the bare allocator fragments. The
underlying cause is that the kernel has no explicit ``retire this thread'' step;
the design for fixing it is recorded in the project notes.

\textbf{Why it rarely bites.} Idiomatic embedded Ada creates a \emph{fixed} set
of tasks at startup and runs them for the life of the system. The fixed-set
case never churns and never leaks. Dynamic create/destroy in a hot loop is an
unusual pattern on a microcontroller; where it is needed, a task pool (allocate
once, reuse) sidesteps the issue entirely.

\section{Minor robustness edges}

\begin{itemize}
  \item \textbf{Lowering the running task's priority below a ready peer} can hang
        the caller (a bare-kernel reschedule edge). Raising priority, or staying
        the highest, works. \emph{Why it rarely bites:} priority is normally set
        once at task creation; the dynamic-lowering-below-peers pattern is rare,
        and priority \emph{ceilings} (the common dynamic case) go up, not down.
  \item \textbf{\texttt{Ada.Text\_IO} over the USB-Serial-JTAG console} is not
        robust for two tasks printing concurrently. \emph{Why it rarely bites:}
        this is a console-transport detail, not the tasking kernel. Route
        diagnostic output through one task, or use a real UART (the UART driver
        is task-safe).
\end{itemize}

\section{Architectural: \texttt{Interrupt\_Priority'Last} and the kernel tick}\index{interrupt priority}

A task declared at \texttt{Interrupt\_Priority'Last} masks the kernel's own
Xtensa level-5 tick and cross-core IPI, deadlocking its activation (ACATS
\texttt{CXD1006}; see the ACATS chapter). The kernel owns the top hardware level
and shares it with the highest Ada interrupt priority.

\textbf{Why it rarely bites.} Only the single \emph{topmost} interrupt priority
collides. Any priority below it (including every priority the peripheral HAL
drivers use for their device interrupts) is unaffected, and ordinary tasks use
\texttt{Priority}, not \texttt{Interrupt\_Priority}, at all. A handler simply
should not claim the one ceiling reserved for the scheduler.

\section{The common case is clean}

None of the above touches the path a typical full-tasking program actually walks:
declare a fixed set of tasks and protected objects at startup; rendezvous and
share data through protected entries; use absolute delays and timing events;
attach interrupt handlers statically; finalize controlled objects at scope exit;
abort at completion points. That path is implemented and hardware-verified. The
limitations are real, but they live at the edges (toolchain corners, an
RM-permitted abort case, and patterns that idiomatic embedded code seldom uses:
library-level finalization, task churn, dynamic priority lowering), and every
one has either a clean workaround or standing permission from the language.
